\documentclass{homework}
\course{Math 6702}
\author{Jim Fowler}
\hwtitle{Differential Geometry}
\input{preamble}

\begin{document}
\maketitle

\begin{inspiration}
  Then it doesn't much matter which way you go.
  \byline{The Cheshire Cat, not exactly talking about the cut locus.}
\end{inspiration}

\section{Terminology}

\begin{problem}
  What are \textbf{normal coordinates}?
\end{problem}

\begin{problem}
  What is a \textbf{minimizing geodesic} between points $p$ and $q$?
\end{problem}

\begin{problem}
  What is the \textbf{cut locus} of a point $p$?
\end{problem}

\section{Numericals}

\begin{problem}
  Consider the hyperbolic plane $\mathbb{H}^2$. Show that for any $p,q \in \mathbb{H}^2$, there is a unique geodesic segment from $p$ to $q$ and that it is minimizing.

  And what does this imply about the cut locus of $\mathbb{H}^2$?
\end{problem}

\begin{problem}
  Let's relate \textbf{length} and \textbf{energy} of curves. Let $(M,g)$ be a Riemannian manifold and let $\gamma:[a,b]\to M$ be smooth.
  Show that
  \[
    L(\gamma)=\int_a^b \|\dot\gamma(t)\|\,dt
    \qquad\text{and}\qquad
    E(\gamma):=\frac12\int_a^b \|\dot\gamma(t)\|^2\,dt
  \]
  satisfy
  \[
    L(\gamma)^2 \le (b-a)\cdot 2E(\gamma),
  \]
  with equality if and only if $\|\dot\gamma(t)\|$ is constant.
\end{problem}

\section{Exploration}

\begin{problem}
  Show that the function $x\mapsto d(m,x)$ is Lipschitz, but also show that $x\mapsto d(m,x)$ may \emph{not} be smooth everywhere.
\end{problem}
  
\begin{problem}
  Suppose $(M,g)$ is a complete Riemannian manifold and $M$ has finite diameter, i.e., there exists $B$ so that for all $x, y \in M$, we have $d(x,y) < B$.

  Does this imply that $M$ is compact?
\end{problem}

\begin{problem}
On a complete Riemannian manifold $(M,g)$, a \textbf{geodesic ray} is an isometric embedding $r:[0,\infty)\to M$.

Define a relation on geodesic rays by $r \sim s$ iff
  \[
    \sup_{t\ge 0} d\bigl(r(t),s(t)\bigr) < \infty.
  \]
  Prove that this is an equivalence relation.
\end{problem}

\begin{problem}
  The \textbf{visual boundary} of $M$ is the set of equivalence classes
  \[
    \partial_\infty M := \{\text{geodesic rays in }M\}/\sim.
  \]
  Fix a basepoint $p\in M$ and let $\mathcal R_p$ denote the set of geodesic rays
  $r:[0,\infty)\to X$ with $r(0)=p$.
  Define 
  \[
    \Phi:\ \mathcal R_p/\sim \ \longrightarrow\ \partial_\infty M,
    \qquad
    [r]\longmapsto \text{the equivalence class of $r$ in $\partial_\infty M$}.
  \]
  Show that $\Phi$ is a bijection.
\end{problem}

\begin{problem}
  Determine $\partial_\infty \mathbb{R}^2$.
\end{problem}

\begin{problem}
  Determine $\partial_\infty \mathbb{H}^2$.
\end{problem}

\begin{problem}
  Determine $\partial_\infty (S^1 \times \mathbb{R})$.
\end{problem}

\begin{problem}
In some cases, the \textbf{marked length spectrum} determines the geometry.

Let $\R^2/\Gamma$ and $\R^2/\Gamma'$ be flat tori, where
$\Gamma,\Gamma'\subset \R^2$ are lattices.
Define for $\gamma\in \Gamma\setminus\{0\}$,
\[
  \ell(\gamma)=\inf\{\text{Length}(\alpha): \alpha \text{ is a loop
  representing the free homotopy class corresponding to }\gamma\}.
\]

Suppose there is a group isomorphism $\phi:\Gamma\to\Gamma'$ such that
\[
  \ell(\phi(\gamma))=\ell(\gamma)\qquad\text{for all }\gamma\in\Gamma\setminus\{0\}.
\]
Show that  $\R^2/\Gamma$ and $\R^2/\Gamma'$ are isometric.
\end{problem}

\section{Prove or Disprove and Salvage if Possible}

\begin{problem}
  In a Riemannian manifold, closed metric balls are compact.
\end{problem}

\begin{problem}
  Let $(M,g)$ be a connected Riemannian manifold and $d$ the distance metric induced from the Riemannian metric $g$.
  Then for every smooth curve $\gamma:[a,b]\to M$,
  \[
    L(\gamma)=d(\gamma(a),\gamma(b)).
  \]
\end{problem}

\end{document}
